\begin{textsample}{POS Dim 2 – pr – Score 41.00 – pr/pr\_em\_3.txt}  \label{ex:f2_pos_111}
A partir das premissas da Base Nacional Comum Curricular ( BRASIL, 2018d ) e do Referencial Curricular para o Ensino Médio do Paraná ( PARANÁ, 2021a ), a Rede Pública Estadual de Ensino elaborou, no segundo \textbf{semestre} de 2021, o \textbf{Currículo} da Formação Geral Básica - FGB ( PARANÁ, 2022a ), com vistas a orientar e apoiar o desenvolvimento da proposta pedagógica curricular das instituições públicas de ensino.
Este \textbf{Currículo}, organizado por áreas do conhecimento, considera o desenvolvimento de competências e habilidades para a formação integral do sujeito, tendo a interdisciplinaridade e a contextualização como princípios \textbf{estruturantes}.
Em cada Área, os componentes curriculares constituem, em seu conjunto, os saberes que, de forma articulada, mobilizam as habilidades.
Nessa organização curricular, os níveis de desenvolvimento das aprendizagens ocorrem nos âmbitos cognitivo, procedimental e atitudinal, sendo expressos na elaboração dos objetivos de aprendizagem e evidenciando o desenvolvimento das habilidades de forma paulatina e progressiva.
O ponto central para a organização da prática pedagógica são, portanto, os objetivos de aprendizagem, sendo os conteúdos, meios para atingi -los ; ambos estão comprometidos com o desenvolvimento das competências e habilidades a serem alcançadas ao longo das três séries do Ensino Médio.
No ano letivo de 2022, conforme a \textbf{Deliberação} CEE/PR n.º 04/2021 ( PARANÁ, 2021 ), o estado do Paraná iniciou a implementação gradativa da nova arquitetura curricular do Ensino Médio.
Na esfera pública, nesta configuração, para a 1ª série do Ensino Médio houve a \textbf{oferta} de componentes curriculares da Formação Geral Básica e de unidades curriculares obrigatórias que integravam os \textbf{Itinerários} \textbf{Formativos} ( denominada Parte \textbf{Flexível} Obrigatória - PFO ), como Projeto de Vida, Educação Financeira e Pensamento Computacional.
As ementas para as unidades curriculares da PFO \textbf{previstas} na Matriz Curricular de cada uma das modalidades de ensino, elaboradas em 2021 a fim de orientar a prática docente, foram organizadas no Caderno de \textbf{Itinerários} \textbf{Formativos} 2022 : Ementas das Unidades Curriculares \textbf{ofertadas} em 2022[^1 ] ( PARANÁ, 2022b ).
Da mesma maneira, ao longo de 2022 e de 2023, para a \textbf{rede} pública \textbf{estadual}, foram elaborados os Cadernos de \textbf{Itinerários} \textbf{Formativos}, respectivamente, para os anos letivos de 2023[^2 ] ( para o trabalho com as turmas da 2ª série ) e 2024[^3 ] ( para as turmas da 3ª série ).
Estes Cadernos apresentaram, além das unidades curriculares da Parte \textbf{Flexível} Obrigatória ( PFO ), as Trilhas de Aprendizagem dos \textbf{Itinerários} \textbf{Formativos} Integrados \textbf{ofertados} pela \textbf{rede} pública de ensino e que eram elegíveis para os estudantes, constituindo -se a Parte \textbf{Flexível} do \textbf{currículo} do Ensino Médio.
As Trilhas de Aprendizagem consideraram as demandas e necessidades da sociedade atual, como também o contexto no qual os estudantes estavam inseridos.
Dessa forma, por meio da indicação de metodologias diversificadas, objetivou -se a ampliação das aprendizagens e a promoção do protagonismo juvenil.
Todo o processo de construção curricular realizado para \textbf{atender} às modalidades de ensino considerou as especificidades de cada público \textbf{atendido} e contou com a participação efetiva das comunidades e povos em sua elaboração.
No ano de 2023, considerando demandas advindas de parte da sociedade, Ministério da Educação ( MEC ), por meio da \textbf{Portaria} n.º 627, de 4 de abril de 2023 ( BRASIL, 2023c ), suspendeu o calendário da implementação do Ensino Médio e realizou consultas públicas, incluindo consultas on-line com estudantes, professores e gestores escolares, reuniões com diferentes \textbf{entidades}, seminários e \textbf{oficinas}.
Essa consulta resultou no Projeto de Lei n.º 5.230/2023 ( BRASIL, 2023a ), que redefiniu a \textbf{política} nacional para o Ensino Médio, ampliando a \textbf{carga} \textbf{horária} dos componentes curriculares obrigatórios, reorganizando o \textbf{currículo} e substituindo os \textbf{Itinerários} \textbf{Formativos} por Percursos de Aprofundamento.
Esse Projeto de Lei foi, então, convertido na Lei n.º 14.945/2024 ( BRASIL, 2024a ).
Com a publicação da Lei Federal n.º 14.945, de 31 de julho de 2024 ( BRASIL, 2024a ), foram \textbf{instituídas} as mudanças para a etapa do Ensino Médio, que impactaram diretamente as arquiteturas curriculares.
A arquitetura curricular é a organização base para a implementação das mudanças previstas na Lei, uma vez que contempla explicitamente o percurso \textbf{formativo} dos estudantes e, implicitamente, a proposta pedagógica a ser desenvolvida.
A fim de orientar essa reorganização curricular dos Sistema de Ensino, foi publicada a Resolução CNE/CEB n.º 2, de 13 de novembro de 2024 ( BRASIL, 2024b ), que \textbf{instituiu} as \textbf{Diretrizes} Curriculares Nacionais para o Ensino Médio ( DCNEM ).
A \textbf{oferta} do Ensino Médio, segundo o Art. 7º da Resolução CNE/CEB n.º 2/2024, assim como toda a Educação Básica, deve ser orientada pelos seguintes princípios gerais : I - a igualdade de condições para acesso, inclusão, permanência e sucesso na escola ; II - a liberdade de aprender, ensinar, pesquisar e divulgar a cultura, o pensamento, a arte e o saber ; III - o pluralismo de ideias e de concepções pedagógicas ; IV - o respeito à liberdade e aos direitos ; V - a coexistência de instituições públicas e privadas de ensino ; VI - a gratuidade do ensino público em estabelecimentos oficiais ; VII - a valorização do \textbf{profissional} da educação escolar ; VIII - a gestão democrática do ensino público, na forma da \textbf{legislação} e normas dos sistemas de ensino ; IX - a \textbf{garantia} de padrão de \textbf{qualidade} ; X - a valorização, na escola, da experiência extraescolar ; e XI - a articulação entre a educação escolar, o trabalho e as práticas sociais. ( BRASIL, 2024b ) É importante, também, destacar os princípios \textbf{orientadores} específicos, associados às singularidades e às necessidades dessa etapa da Educação Básica, presentes no Art. 8º : I. a Formação Integral e Integrada dos estudantes, assegurando a articulação e a integração entre a FGB e os \textbf{itinerários} \textbf{formativos}, a interdisciplinaridade e a contextualização ; II - a indissociabilidade entre educação e prática social, considerando a historicidade dos conhecimentos e dos sujeitos do processo educativo, bem como entre teoria e prática no processo de ensino-aprendizagem ; III - o reconhecimento, valorização e mobilização permanente e integrada das dimensões \textbf{formativas} próprias do mundo do trabalho, na ciência, na tecnologia e na cultura ; IV - a justiça curricular e busca permanente da \textbf{equidade} educacional ; V - o reconhecimento e valorização das diferenças e da diversidade dos sujeitos da ação educativa ; nas múltiplas dimensões de suas identidades, experiências e singularidades ; VI - a afirmação, valorização e defesa da democracia e da cultura de promoção dos direitos humanos ; VII - a \textbf{garantia} de padrões adequados de aprendizagem e desenvolvimento para todos os educandos e todas as educandas ; VIII - a \textbf{garantia} de processos de transição dos anos finais do Ensino Fundamental para o Ensino Médio, considerando as necessidades, as singularidades e as especificidades dos educandos e educandas ; IX - a \textbf{integralidade} e visão sistêmica da proteção às \textbf{trajetórias} escolares no Ensino Médio, com \textbf{garantia} de ações para a permanência, aprendizagem e conclusão do Ensino Médio na idade adequada ; X - o trabalho como princípio educativo e pesquisa como princípio pedagógico ; e XI - a indissociabilidade das dimensões do trabalho, da cultura, da ciência e da tecnologia na formação dos educandos, considerando : a ) o trabalho, em todas as suas formas de organização e expressão, na perspectiva ontológica de transformação da natureza e criação da cultura, como realização inerente ao ser humano e como mediação no processo de produção da sua existência sócio-histórica ; b ) a ciência como o conjunto de conhecimentos sistematizados, produzidos socialmente ao longo da história, na busca da compreensão e da transformação da natureza e da sociedade ; c ) a tecnologia como a transformação da ciência em força produtiva ou mediação do conhecimento científico e a produção, marcada desde sua origem, pelas relações sociais que a levaram a ser produzida ; e d ) a cultura como o processo de produção de expressões materiais, símbolos, representações e significados que correspondem a valores éticos, políticos e estéticos que orientam as normas de conduta de uma sociedade. ( BRASIL, 2024b ) Partindo desses princípios, a organização curricular do Ensino Médio deve ser estruturada tendo -se em vista a Formação Integral dos estudantes.
Dessa maneira, são necessárias a articulação e a integração entre a Formação Geral Básica e os \textbf{Itinerários} \textbf{Formativos}, tanto de Aprofundamento quanto de Formação \textbf{Técnica} e Profissional.
Assim, a \textbf{oferta} da FGB e dos \textbf{Itinerários} não pode ocorrer de maneira segregada, em blocos distintos ( BRASIL, 2024b ).
Essa \textbf{oferta} também precisa assegurar os direitos e objetivos de aprendizagem presentes na BNCC a todos os estudantes, considerando -se todas as formas de \textbf{oferta} e modalidades do Ensino Médio, e estando em acordo com as \textbf{diretrizes} curriculares nacionais de cada uma delas, ressaltando -se o princípio da justiça curricular, como considerado no Art. 5º da Resolução CNE/CEB n.º 2/2024 : IV - justiça curricular : princípio de organização do \textbf{currículo} que estabelece como parâmetros para a tomada de decisões da gestão educacional, da gestão escolar e das práticas pedagógicas : a ) a priorização de conhecimentos e conteúdos de ensino orientados para a promoção, defesa e compromisso com a \textbf{garantia} de uma vida digna para todas as pessoas ; b ) a explicitação e a materialização de uma ética do cuidado e do bem viver nas relações entre o estado e a sociedade ; e c ) a construção de uma convivência solidária e democrática, comprometida com a realização cotidiana dos direitos humanos e a superação das múltiplas formas de exclusão, discriminação, preconceitos e opressão ( BRASIL, 2024b ) As propostas curriculares dos Sistemas de Ensino também devem “ assegurar os direitos de aprendizagem por meio da progressão adequada das competências e habilidades das diferentes áreas do conhecimento ”, conforme o Art. 11, além de garantir que sejam observados e promovidos : o protagonismo dos estudantes ; a mobilização, orientação e apoio na construção de seus projetos de vida ; o tratamento interdisciplinar nas/das diferentes áreas do conhecimento ; a mobilização dos temas contemporâneos transversais estabelecidos na BNCC ; a avaliação da aprendizagem por meio de metodologias diversificadas, de caráter \textbf{formativo} ou somativo e que considerem a singularidade dos sujeitos nas diversidades ; e a ampliação e expansão dos espaços e tempos escolares.
A oportunidade de construção do projeto de vida pelo estudante também deve ser garantida pelos Sistemas de Ensino, sendo necessário “ promover processos intencionais e estruturados de aprendizagem e desenvolvimento integral dos jovens ” ( BRASIL, 2024b ), que desenvolvam a capacidade de dar sentido à sua existência, fazer reflexões coletivas e individuais acerca das questões que permeiam a sua realidade, tomar decisões, planejar o futuro e agir no presente com autonomia e responsabilidade.
A Formação Geral Básica ( FGB ) teve sua \textbf{carga} \textbf{horária} mínima ampliada para 2.400 ( duas mil e quatrocentas ) horas, implicando em uma modificação substancial nas Matrizes Curriculares.
Os direitos e objetivos de aprendizagem expressos nas competências e habilidades da BNCC devem ser desenvolvidos considerando as 4 ( quatro ) áreas do conhecimento, tendo em conta o tratamento interdisciplinar em seus componentes curriculares obrigatórios, respectivamente destacados no Art. 17 da Resolução CNE/CEB n.º 2/2024 : I - linguagens e suas tecnologias, integrada pelos componentes curriculares obrigatórios de língua portuguesa e suas literaturas, língua inglesa, artes e educação física ; II - matemática e suas tecnologias, com o componente curricular obrigatório de matemática ; III - ciências da natureza e suas tecnologias, integrada pelos componentes curriculares obrigatórios de biologia, física e química ; e IV - ciências humanas e sociais aplicadas, integrada pelos componentes curriculares obrigatórios de filosofia, geografia, história e sociologia. ( BRASIL, 2024b ) É importante destacar que a FGB também deve focar no desenvolvimento de competências digitais, que envolvem o letramento digital, o ensino da computação, programação, robótica entre outras, de acordo com a Lei n.º 14.533, de 11 de janeiro de 2023, que \textbf{instituiu} a Política Nacional de Educação Digital ( BRASIL, 2023b ).
Os \textbf{Itinerários} \textbf{Formativos} de Aprofundamento devem ser articulados e integrados à Formação Geral Básica, os \textbf{Itinerários} de Formação \textbf{Técnica} e \textbf{Profissional} também deverão buscar essa coesão curricular e pedagógica com a FGB, a fim de garantir a formação integral dos estudantes para o exercício da cidadania, progressão dos estudos e preparação para o mundo do trabalho.
Além disso, a Resolução CNE/CEB n.º 2/2024 ( BRASIL, 2024b ) ressalta que a \textbf{oferta} dos \textbf{Itinerários} de Formação \textbf{Técnica} e \textbf{Profissional} para estudantes indígenas, quilombolas e do campo deve observar as \textbf{Diretrizes} Curriculares Nacionais dessas modalidades.
Da mesma forma, é fundamental considerar as particularidades dos estudantes da educação especial na perspectiva da educação inclusiva e da Educação Bilíngue de Surdos, respeitando suas necessidades e especificidades.
Ainda sobre a \textbf{oferta}, o Ensino Médio regular, considerando a indissociabilidade entre a FGB e os \textbf{Itinerários} \textbf{Formativos}, terá duração mínima de 3 ( três ) anos, com “ \textbf{carga} \textbf{horária} mínima total de 3.000 ( três mil ) horas, tendo como referência uma \textbf{carga} \textbf{horária} anual de 1.000 ( mil ) horas, distribuídas em pelo menos 200 ( \textbf{duzentos} ) dias de efetivo trabalho escolar ” ( BRASIL, 2024b ).
Essa \textbf{carga} \textbf{horária} deve ser garantida igualmente ao Ensino Médio noturno, podendo ser \textbf{ofertada} em mais de 3 ( três ) anos ; contudo, para \textbf{atender} com \textbf{qualidade} as singularidades e condições dos estudantes \textbf{atendidos} neste período é necessária uma diferenciação curricular e metodológica.
Já o Ensino Médio regular diurno pode ser \textbf{ofertado} também em \textbf{regime} de tempo integral, com, no mínimo, 7 ( sete ) horas diárias, quando adequado aos seus estudantes.
O Plano Nacional de Educação para o próximo decênio estabelecerá metas para essa ampliação progressiva da \textbf{carga} \textbf{horária} anual total do Ensino Médio para 1.400 ( mil e quatrocentas ) horas.
Na Resolução CNE/CEB n.º 2/2024 ( BRASIL, 2024b ) reafirma -se a priorização da \textbf{oferta} regular e presencial do Ensino Médio a todos os estudantes \textbf{adolescentes}, jovens e adultos de todas as modalidades e formas de \textbf{oferta} do ensino.
Porém, em seu Art. 28, § 6º, inciso V, afirma -se que : > em situações excepcionais, respeitados os parâmetros legais \textbf{vigentes} no país e as \textbf{diretrizes} curriculares específicas das diferentes modalidades da Educação Básica, a educação mediada por tecnologia pode ser utilizada para assegurar o direito à educação em regiões de difícil acesso, para o \textbf{currículo} do Ensino Médio na modalidade EJA. ( BRASIL, 2024b ) Na elaboração de suas propostas pedagógicas, as unidades escolares devem considerar a proposta curricular definida pelo Sistema de Ensino, garantindo sua adequação às normas curriculares nacionais e às \textbf{diretrizes} do sistema educacional de seu território.
Esse processo deve respeitar o princípio do pluralismo de ideias e concepções pedagógicas, assegurando o exercício da autonomia e a gestão democrática.
Além disso, é importante que a construção dessas propostas esteja \textbf{alinhada} às demandas da comunidade escolar, promovendo uma educação contextualizada e participativa.
A Resolução CNE/CEB n.º 2/2024 ( BRASIL, 2024b ) ainda discorre sobre os processos de avaliação educacional e da aprendizagem.
Em um processo avaliativo são coletadas evidências da eficiência das ações, tendo em vista uma tomada de decisão.
Assim, a avaliação da aprendizagem e do desenvolvimento dos estudantes deve ser um processo \textbf{formativo} e somativo, utilizando -se diferentes instrumentos e métodos, contemplando as suas características, singularidades e necessidades.
É necessário o registro das evidências, garantindo a \textbf{documentação} pedagógica do processo de ensino e aprendizagem para o acompanhamento adequado dos estudantes, seja para a verificação dos seus avanços e necessidades ao longo do ano letivo, seja para a tomada de decisões a respeito da sua progressão ou da necessidade de estratégias específicas de apoio complementar.
O processo de avaliação educacional se estende para além da sala de aula, contemplando avaliações institucionais e participativas da escola, visando a melhoria contínua da organização, funcionamento e resultados educacionais alcançados, e avaliações externas, em larga escala, para uma avaliação das \textbf{Redes} e Sistemas de Ensino, a fim de embasar as \textbf{políticas} educacionais em todas as esferas ( \textbf{municipais}, \textbf{estaduais} e nacionais, de gestão pública ou privada ).
Os processos avaliativos, em conjunto com estratégias permanentes de monitoramento de dados educacionais, também auxiliam os Sistemas de Ensino a estabelecer programas e ações de acompanhamento, a fim de garantir a permanência estudantil e prevenir o abandono, a evasão e a \textbf{reprovação} no Ensino Médio.
Ações, estratégias e programas voltados à busca ativa, à identificação e intervenção precoce dos estudantes em risco de evasão e abandono, ao acompanhamento do acesso, da permanência e da superação da retenção escolar são reforçadas na Resolução CNE/CEB n.º 2/2024 ( BRASIL, 2024b ) e devem \textbf{atender} à democratização do acesso, permanência e sucesso escolar com \textbf{qualidade} social da educação.
Para o Paraná, considerando esses aspectos e as características do território, foram estruturados demarcadores conceituais que fundamentam a Política Educacional para o Ensino Médio do estado do Paraná.
Os demarcadores conceituais permitem o desenvolvimento de uma prática pedagógica intencional, já que explicitam a compreensão do estado sobre a sociedade, a aprendizagem e o protagonismo juvenil, a organização curricular, os sujeitos do Ensino Médio e a \textbf{qualidade} social da educação.
Estes demarcadores serão explorados no segundo \textbf{capítulo} deste Referencial Curricular. [ ^1 ] : O Caderno de \textbf{Itinerários} \textbf{Formativos} 2022 : Ementas das Unidades Curriculares \textbf{ofertadas} em 2022 foi \textbf{atualizado} e passou a integrar os Cadernos de 2023. [ ^2 ] : Os Cadernos de \textbf{Itinerários} \textbf{Formativos} 2023 : Ementas das Unidades Curriculares estão disponíveis em : https://acervodigital.educacao.pr.gov.br/pages/search.php? k=b0c4d14f9b\&search=\%21collection3701. ( PARANÁ, 2023a ) [ ^3 ] : Os Cadernos de \textbf{Itinerários} \textbf{Formativos} 2024 : Ementas das Unidades Curriculares estão disponíveis em : https://acervodigital.educacao.pr.gov.br/pages/search.php? k=b0c4d14f9b\&search=\%21collection6613. ( PARANÁ, 2024 )

% matched lemmas: adolescente, alinhar, atender, atualizar, capítulo, carga, currículo, deliberação, diretriz, documentação, duzentos, entidade, equidade, estadual, estruturante, flexível, formativo, garantia, horário, instituir, integralidade, itinerário, legislação, municipal, oferta, ofertar, oficina, orientador, política, portaria, previsto, profissional, qualidade, rede, regime, reprovação, semestre, trajetória, técnico, vigente
\end{textsample}
